% Figure 04 — Diagramme de cas d'utilisation UML
% Acteurs : Visiteur, Utilisateur authentifié
\documentclass[border=10pt]{standalone}
% Préambule commun pour tous les diagrammes TikZ standalone
% À inclure via % Préambule commun pour tous les diagrammes TikZ standalone
% À inclure via % Préambule commun pour tous les diagrammes TikZ standalone
% À inclure via \input{preambule.tex} après \documentclass{standalone}
\usepackage[utf8]{inputenc}
\usepackage[T1]{fontenc}
\usepackage[french]{babel}
\usepackage{lmodern}
\usepackage{tikz}
\usepackage{xcolor}

\usetikzlibrary{
  arrows.meta,
  positioning,
  shapes.geometric,
  shapes.symbols,
  shapes.misc,
  fit,
  calc,
  backgrounds,
  decorations.pathreplacing,
  decorations.markings,
  shadows.blur,
  patterns
}

% Palette de couleurs (cohérente sur tout le rapport)
\definecolor{awsorange}{HTML}{FF9900}
\definecolor{awsdark}{HTML}{232F3E}
\definecolor{awsblue}{HTML}{1A476F}
\definecolor{dockerblue}{HTML}{2496ED}
\definecolor{dockerdark}{HTML}{0DB7ED}
\definecolor{nginxgreen}{HTML}{009639}
\definecolor{flaskgray}{HTML}{3776AB}
\definecolor{pgblue}{HTML}{336791}
\definecolor{tlsgreen}{HTML}{2E8B57}
\definecolor{netgray}{HTML}{6C757D}
\definecolor{bgsoft}{HTML}{F5F7FA}
\definecolor{bordersoft}{HTML}{C8D1DC}
\definecolor{textdark}{HTML}{1F2937}
\definecolor{accentred}{HTML}{C0392B}
\definecolor{accentyellow}{HTML}{F1C40F}

% Styles génériques réutilisés dans plusieurs diagrammes
\tikzset{
  base/.style={
    rounded corners=4pt,
    draw=bordersoft,
    fill=white,
    text=textdark,
    align=center,
    inner sep=6pt,
    font=\small
  },
  zone/.style={
    base,
    fill=bgsoft,
    draw=netgray,
    dashed,
    inner sep=12pt
  },
  cloud/.style={
    base,
    fill=awsdark!5,
    draw=awsorange,
    line width=0.8pt,
    inner sep=14pt
  },
  service/.style={
    base,
    line width=0.6pt,
    minimum width=2.6cm,
    minimum height=1.1cm,
    drop shadow={shadow xshift=1pt, shadow yshift=-1pt, opacity=0.2}
  },
  legendebox/.style={
    base,
    fill=white,
    font=\footnotesize,
    inner sep=4pt
  },
  fleche/.style={
    -{Latex[length=2.5mm,width=2mm]},
    line width=0.6pt,
    color=netgray
  },
  fleche-bidir/.style={
    {Latex[length=2.5mm,width=2mm]}-{Latex[length=2.5mm,width=2mm]},
    line width=0.6pt,
    color=netgray
  },
  fleche-bold/.style={
    -{Latex[length=3mm,width=2.4mm]},
    line width=1pt
  },
  etiquette/.style={
    font=\scriptsize\sffamily,
    fill=white,
    inner sep=2pt,
    text=textdark
  }
}
 après \documentclass{standalone}
\usepackage[utf8]{inputenc}
\usepackage[T1]{fontenc}
\usepackage[french]{babel}
\usepackage{lmodern}
\usepackage{tikz}
\usepackage{xcolor}

\usetikzlibrary{
  arrows.meta,
  positioning,
  shapes.geometric,
  shapes.symbols,
  shapes.misc,
  fit,
  calc,
  backgrounds,
  decorations.pathreplacing,
  decorations.markings,
  shadows.blur,
  patterns
}

% Palette de couleurs (cohérente sur tout le rapport)
\definecolor{awsorange}{HTML}{FF9900}
\definecolor{awsdark}{HTML}{232F3E}
\definecolor{awsblue}{HTML}{1A476F}
\definecolor{dockerblue}{HTML}{2496ED}
\definecolor{dockerdark}{HTML}{0DB7ED}
\definecolor{nginxgreen}{HTML}{009639}
\definecolor{flaskgray}{HTML}{3776AB}
\definecolor{pgblue}{HTML}{336791}
\definecolor{tlsgreen}{HTML}{2E8B57}
\definecolor{netgray}{HTML}{6C757D}
\definecolor{bgsoft}{HTML}{F5F7FA}
\definecolor{bordersoft}{HTML}{C8D1DC}
\definecolor{textdark}{HTML}{1F2937}
\definecolor{accentred}{HTML}{C0392B}
\definecolor{accentyellow}{HTML}{F1C40F}

% Styles génériques réutilisés dans plusieurs diagrammes
\tikzset{
  base/.style={
    rounded corners=4pt,
    draw=bordersoft,
    fill=white,
    text=textdark,
    align=center,
    inner sep=6pt,
    font=\small
  },
  zone/.style={
    base,
    fill=bgsoft,
    draw=netgray,
    dashed,
    inner sep=12pt
  },
  cloud/.style={
    base,
    fill=awsdark!5,
    draw=awsorange,
    line width=0.8pt,
    inner sep=14pt
  },
  service/.style={
    base,
    line width=0.6pt,
    minimum width=2.6cm,
    minimum height=1.1cm,
    drop shadow={shadow xshift=1pt, shadow yshift=-1pt, opacity=0.2}
  },
  legendebox/.style={
    base,
    fill=white,
    font=\footnotesize,
    inner sep=4pt
  },
  fleche/.style={
    -{Latex[length=2.5mm,width=2mm]},
    line width=0.6pt,
    color=netgray
  },
  fleche-bidir/.style={
    {Latex[length=2.5mm,width=2mm]}-{Latex[length=2.5mm,width=2mm]},
    line width=0.6pt,
    color=netgray
  },
  fleche-bold/.style={
    -{Latex[length=3mm,width=2.4mm]},
    line width=1pt
  },
  etiquette/.style={
    font=\scriptsize\sffamily,
    fill=white,
    inner sep=2pt,
    text=textdark
  }
}
 après \documentclass{standalone}
\usepackage[utf8]{inputenc}
\usepackage[T1]{fontenc}
\usepackage[french]{babel}
\usepackage{lmodern}
\usepackage{tikz}
\usepackage{xcolor}

\usetikzlibrary{
  arrows.meta,
  positioning,
  shapes.geometric,
  shapes.symbols,
  shapes.misc,
  fit,
  calc,
  backgrounds,
  decorations.pathreplacing,
  decorations.markings,
  shadows.blur,
  patterns
}

% Palette de couleurs (cohérente sur tout le rapport)
\definecolor{awsorange}{HTML}{FF9900}
\definecolor{awsdark}{HTML}{232F3E}
\definecolor{awsblue}{HTML}{1A476F}
\definecolor{dockerblue}{HTML}{2496ED}
\definecolor{dockerdark}{HTML}{0DB7ED}
\definecolor{nginxgreen}{HTML}{009639}
\definecolor{flaskgray}{HTML}{3776AB}
\definecolor{pgblue}{HTML}{336791}
\definecolor{tlsgreen}{HTML}{2E8B57}
\definecolor{netgray}{HTML}{6C757D}
\definecolor{bgsoft}{HTML}{F5F7FA}
\definecolor{bordersoft}{HTML}{C8D1DC}
\definecolor{textdark}{HTML}{1F2937}
\definecolor{accentred}{HTML}{C0392B}
\definecolor{accentyellow}{HTML}{F1C40F}

% Styles génériques réutilisés dans plusieurs diagrammes
\tikzset{
  base/.style={
    rounded corners=4pt,
    draw=bordersoft,
    fill=white,
    text=textdark,
    align=center,
    inner sep=6pt,
    font=\small
  },
  zone/.style={
    base,
    fill=bgsoft,
    draw=netgray,
    dashed,
    inner sep=12pt
  },
  cloud/.style={
    base,
    fill=awsdark!5,
    draw=awsorange,
    line width=0.8pt,
    inner sep=14pt
  },
  service/.style={
    base,
    line width=0.6pt,
    minimum width=2.6cm,
    minimum height=1.1cm,
    drop shadow={shadow xshift=1pt, shadow yshift=-1pt, opacity=0.2}
  },
  legendebox/.style={
    base,
    fill=white,
    font=\footnotesize,
    inner sep=4pt
  },
  fleche/.style={
    -{Latex[length=2.5mm,width=2mm]},
    line width=0.6pt,
    color=netgray
  },
  fleche-bidir/.style={
    {Latex[length=2.5mm,width=2mm]}-{Latex[length=2.5mm,width=2mm]},
    line width=0.6pt,
    color=netgray
  },
  fleche-bold/.style={
    -{Latex[length=3mm,width=2.4mm]},
    line width=1pt
  },
  etiquette/.style={
    font=\scriptsize\sffamily,
    fill=white,
    inner sep=2pt,
    text=textdark
  }
}


\usetikzlibrary{shapes.geometric}

\begin{document}
\begin{tikzpicture}[node distance=6mm and 12mm,
  acteur/.style={
    draw=textdark, fill=white, line width=0.8pt,
    minimum width=1.4cm, minimum height=1.6cm,
    align=center, font=\small
  },
  uc/.style={
    ellipse, draw=awsblue, fill=awsblue!8,
    line width=0.6pt, align=center,
    minimum width=3.4cm, minimum height=0.9cm,
    font=\footnotesize, inner sep=2pt
  },
  systeme/.style={
    draw=netgray, line width=0.8pt, dashed,
    rounded corners=4pt, fill=bgsoft,
    inner sep=14pt
  }
]

  % --- Acteur Visiteur ---
  \node (visitor-head) [circle, draw, minimum size=0.7cm] {};
  \node[below=0pt of visitor-head, draw, minimum width=0.6cm,
        minimum height=0.9cm] (visitor-body) {};
  \node[below=0pt of visitor-body, font=\small\bfseries] (visitor-label) {Visiteur};
  \node[fit=(visitor-head)(visitor-label),
        inner sep=1pt] (visitor) {};

  % --- Acteur Utilisateur authentifié ---
  \node[right=12cm of visitor-head] (user-head)
        [circle, draw, minimum size=0.7cm] {};
  \node[below=0pt of user-head, draw, minimum width=0.6cm,
        minimum height=0.9cm] (user-body) {};
  \node[below=0pt of user-body, font=\small\bfseries, align=center]
        (user-label) {Utilisateur\\authentifié};
  \node[fit=(user-head)(user-label),
        inner sep=1pt] (user) {};

  % --- Cas d'utilisation (sans frame, on les ajoute ensuite) ---
  \node[uc] (uc-register) at (3,1.5) {S'inscrire};
  \node[uc] (uc-login)    at (3,0)   {Se connecter};
  \node[uc] (uc-logout)   at (9,3)   {Se déconnecter};

  \node[uc] (uc-create)  at (6,1.6)  {Créer une note};
  \node[uc] (uc-read)    at (6,0.4)  {Consulter ses notes};
  \node[uc] (uc-edit)    at (6,-0.8) {Éditer une note};
  \node[uc] (uc-delete)  at (6,-2.0) {Supprimer une note};

  \node[uc] (uc-archive) at (9,-0.4) {Archiver / désarchiver};
  \node[uc] (uc-trash)   at (9,-1.6) {Corbeille \(+\) restauration};
  \node[uc] (uc-search)  at (9,1.6)  {Rechercher une note};

  % --- Frame système ---
  \node[systeme, fit=(uc-register)(uc-login)(uc-logout)
                     (uc-create)(uc-read)(uc-edit)(uc-delete)
                     (uc-archive)(uc-trash)(uc-search),
        label={[font=\small\bfseries, text=netgray]north:Système — WebCyber Notes}
       ] (sys) {};

  % --- Liens acteur Visiteur ---
  \draw[fleche, color=textdark] (visitor.east) -- (uc-register.west);
  \draw[fleche, color=textdark] (visitor.east) -- (uc-login.west);

  % --- Liens acteur Utilisateur ---
  \draw[fleche, color=textdark] (user.west) -- (uc-logout.east);
  \draw[fleche, color=textdark] (user.west) -- (uc-create.east);
  \draw[fleche, color=textdark] (user.west) -- (uc-read.east);
  \draw[fleche, color=textdark] (user.west) -- (uc-edit.east);
  \draw[fleche, color=textdark] (user.west) -- (uc-delete.east);
  \draw[fleche, color=textdark] (user.west) -- (uc-archive.east);
  \draw[fleche, color=textdark] (user.west) -- (uc-trash.east);
  \draw[fleche, color=textdark] (user.west) -- (uc-search.east);

  % --- Relation include : Login << include >> Auth check ---
  \draw[-{Latex[length=2mm]}, dashed, color=accentred]
    (uc-create) -- node[etiquette, sloped, font=\tiny, text=accentred]
                 {\guillemotleft\,include\,\guillemotright}
                 (uc-login);
  \draw[-{Latex[length=2mm]}, dashed, color=accentred]
    (uc-edit) -- node[etiquette, sloped, font=\tiny, text=accentred]
                 {\guillemotleft\,include\,\guillemotright}
                 (uc-login);

\end{tikzpicture}
\end{document}
